\documentclass[10pt]{article}
\providecommand{\FIGDIR}{fig}
\newcommand{\DIGESTDATE}{23 August 2026}
% ===================== PM Research Watch - shared preamble =====================
\usepackage[T1]{fontenc}
\usepackage[utf8]{inputenc}
\usepackage{amsmath,amssymb}
\usepackage{tgpagella}
\usepackage{tgheros}
\renewcommand{\sfdefault}{qhv}
\renewcommand{\ttdefault}{lmtt}
\usepackage{microtype}
\usepackage[english]{babel}

\usepackage[a4paper,top=17mm,bottom=17mm,left=17mm,right=17mm,
            headsep=6mm,footskip=10mm]{geometry}
\usepackage{graphicx}
\usepackage{xcolor}
\usepackage{colortbl}
\usepackage{tikz}
\usetikzlibrary{calc}
\usepackage[most]{tcolorbox}
\usepackage{titlesec}
\usepackage{fancyhdr}
\usepackage{booktabs}
\usepackage{array}
\usepackage{longtable}
\usepackage{enumitem}
\usepackage{multicol}
\usepackage{ragged2e}
\usepackage[hidelinks]{hyperref}

% ------------------------------------------------------------------ palette
\definecolor{Ink}    {HTML}{14262E}
\definecolor{Deep}   {HTML}{0F4C5C}
\definecolor{Teal}   {HTML}{1F7A8C}
\definecolor{Sky}    {HTML}{4FA3C4}
\definecolor{Sage}   {HTML}{6FA287}
\definecolor{Amber}  {HTML}{E3A008}
\definecolor{Clay}   {HTML}{B5643C}
\definecolor{Coral}  {HTML}{D1495B}
\definecolor{Violet} {HTML}{6A5B9E}
\definecolor{Slate}  {HTML}{4A6572}
\definecolor{Paper}  {HTML}{FBF9F5}
\definecolor{Mist}   {HTML}{EEEAE2}
\definecolor{Rule}   {HTML}{D9D3C9}

\pagecolor{Paper}
\color{Ink}

\hypersetup{colorlinks=true, linkcolor=Deep, urlcolor=Teal, citecolor=Teal,
            pdftitle={PM Research Watch}, pdfauthor={Automated daily digest}}

% ------------------------------------------------------------------ headings
\titleformat{\section}
  {\sffamily\bfseries\large\color{Deep}}{}{0pt}
  {}[{\vspace{1.2mm}\color{Amber}\titlerule[1.1pt]}]
\titlespacing*{\section}{0pt}{5.5mm}{2.8mm}

\titleformat{\subsection}
  {\sffamily\bfseries\normalsize\color{Teal}}{}{0pt}{}
\titlespacing*{\subsection}{0pt}{3.6mm}{1.4mm}

\setlist[itemize]{leftmargin=4.2mm, itemsep=1.1mm, topsep=1.3mm, parsep=0pt,
                  label=\textcolor{Teal}{\footnotesize$\blacktriangleright$}}

% ------------------------------------------------------------------ header/footer
\pagestyle{fancy}
\fancyhf{}
\renewcommand{\headrulewidth}{0pt}
\renewcommand{\footrulewidth}{0pt}
\fancyhead[L]{\sffamily\scriptsize\color{Slate}PM RESEARCH WATCH\;\textcolor{Rule}{\textbar}\;\DIGESTDATE}
\fancyhead[R]{\sffamily\scriptsize\color{Slate}Particulate matter monitoring \& health effects}
\fancyfoot[C]{\sffamily\scriptsize\color{Slate}\thepage}

% ------------------------------------------------------------------ components
% metric tile
\newcommand{\metric}[3]{%
  \begin{tcolorbox}[enhanced, colback=white, colframe=Rule, boxrule=0.5pt,
    arc=1.6mm, left=2.2mm, right=2.2mm, top=1.8mm, bottom=1.8mm,
    borderline west={1.7mm}{0pt}{#3}, nobeforeafter, valign=center,
    height=17mm, width=\linewidth]
    \RaggedRight
    {\sffamily\bfseries\fontsize{17}{17}\selectfont\color{#3}#1}\\[0.5mm]
    {\sffamily\fontsize{6.5}{7.4}\selectfont\color{Slate}#2}
  \end{tcolorbox}}

% pull-quote / key insight box
\newtcolorbox{keybox}[1][]{enhanced, colback=Mist, colframe=Mist,
  arc=1.6mm, left=3.2mm, right=3.2mm, top=2.4mm, bottom=2.4mm,
  borderline west={1.7mm}{0pt}{Amber}, #1}

% subtopic band
\newcommand{\band}[2]{%
  \vspace{3.2mm}\noindent
  \begin{tcolorbox}[enhanced, colback=#2, colframe=#2, arc=1.2mm,
    left=2.6mm, right=2.6mm, top=1.1mm, bottom=1.1mm, boxrule=0pt,
    width=\linewidth, nobeforeafter]
    {\sffamily\bfseries\small\color{white}#1}
  \end{tcolorbox}\vspace{1.6mm}}

% one paper entry: {short cite}{title}{journal}{doi}{body}
\newcommand{\paperentry}[5]{%
  \par\vspace{1.9mm}\noindent
  {\sffamily\bfseries\footnotesize\color{Ink}#1}\;%
  {\sffamily\fontsize{7}{8}\selectfont\color{Slate}#3}\\[0.4mm]
  {\footnotesize\itshape\color{Deep}#2}\\[0.6mm]
  {\footnotesize #4}\\[0.3mm]
  {\fontsize{6.6}{7.6}\selectfont\color{Slate}\texttt{doi:}\,\href{https://doi.org/#5}{\texttt{\detokenize{#5}}}}
  \par}

% caption style
\newcommand{\figcap}[1]{\par\vspace{1.0mm}
  {\fontsize{7.2}{8.6}\selectfont\color{Slate}\textsf{\textbf{Fig.}}\; #1\par}\vspace{1.5mm}}

\providecommand{\permil}{\textperthousand}

\setlength{\parindent}{0pt}
\setlength{\parskip}{0pt}
\renewcommand{\arraystretch}{1.18}


\begin{document}

% ==================================================================== MASTHEAD
\thispagestyle{empty}
\begin{tikzpicture}[remember picture, overlay]
  \fill[Deep] (current page.north west) rectangle
      ($(current page.north east)+(0,-67mm)$);
  \fill[Amber] ($(current page.north west)+(0,-67mm)$) rectangle
      ($(current page.north east)+(0,-68.4mm)$);
  \foreach \x/\y/\r in {22/12/0.9, 41/26/1.5, 63/9/0.7, 88/31/1.1, 112/17/1.9,
                        138/29/0.8, 159/13/1.3, 178/54/1.0, 196/21/1.6,
                        30/44/0.6, 74/57/1.2, 125/48/0.9, 168/6/0.8, 191/59/0.7,
                        52/18/1.0, 100/60/0.6, 148/20/1.1, 66/49/1.4, 205/45/0.8}
    \fill[Sky, opacity=0.30] ($(current page.north west)+(\x mm,-\y mm)$) circle (\r mm);
\end{tikzpicture}

\vspace*{4mm}
{\sffamily\fontsize{7.6}{9}\selectfont\color{Sky}%
 AUTOMATED DAILY LITERATURE DIGEST\;\textbullet\;PhD RESEARCH WATCH\par}
\vspace{2.2mm}
{\sffamily\bfseries\fontsize{27}{29}\selectfont\color{white}%
 Particulate Matter\par}
{\sffamily\fontsize{27}{29}\selectfont\color{Sky}%
 Monitoring \& Health Effects\par}
\vspace{5.0mm}
{\sffamily\fontsize{8.4}{10}\selectfont\color{white}%
 \mbox{Issue 2026-08-23}\;\;\textcolor{Amber}{\textbar}\;\;\mbox{Entry window 23 Aug 2026}\;\;
 \textcolor{Amber}{\textbar}\;\;\mbox{9 in scope, 22 screened out}\par}
\vspace{18mm}

% -------------------------------------------------------------- metric strip
\noindent
\begin{minipage}[t]{0.190\linewidth}\metric{9}{RECORDS IN SCOPE\\(22 SCREENED OUT)}{Deep}\end{minipage}\hfill
\begin{minipage}[t]{0.190\linewidth}\metric{4}{OF 10 CLUSTERS\\OCCUPIED}{Teal}\end{minipage}\hfill
\begin{minipage}[t]{0.190\linewidth}\metric{2}{ABSTRACT-FREE\\DEPOSITS (22\%)}{Sky}\end{minipage}\hfill
\begin{minipage}[t]{0.190\linewidth}\metric{0}{RISK ESTIMATES\\ON A PM CONTRAST}{Amber}\end{minipage}\hfill
\begin{minipage}[t]{0.190\linewidth}\metric{3}{INSTRUMENT- OR\\METRIC-LEVEL PAPERS}{Clay}\end{minipage}

\vspace{6mm}

% -------------------------------------------------------------- signal of day
\section*{\sffamily\bfseries\color{Deep}Signal of the day}
\vspace{-1.5mm}{\color{Amber}\rule{\linewidth}{1.1pt}}\vspace{2.5mm}

\begin{keybox}
{\sffamily\bfseries\small\color{Deep}A Sunday issue with no forest plot and an unusually
direct message for the sensing side: the quantity being regulated and the quantity that
drives inflammation are coming apart, and today's instrument paper is measuring the wrong
one.}\\[1.4mm]
\footnotesize
Zhong et al.\ (\textit{Environ Int}) ran three years of Fuzhou observations against an
animal-derived surface-area/inflammation dose--response and found that \textbf{6.6\% of
hours that satisfied the PM$_{2.5}$ standard ($<$35\,$\mu$g/m$^3$) carried
lung-deposited surface area 7.5$\times$ above what the local LDSA/PM$_{2.5}$ ratio
predicts}, so that \textbf{82\% of the predicted inflammatory burden was invisible to
mass-based assessment}. In those hours median PM$_{2.5}$ was \textbf{3.0\,$\mu$g/m$^3$}
--- an order of magnitude under the standard --- while the dose--response projects 27\%
lung neutrophil infiltration. Their ML attribution puts the mechanism on summer--autumn
afternoon photochemistry producing 50--100\,nm particles, i.e.\ a regime that
\emph{grows} as primary emissions fall and radiation load rises. Every optical low-cost
sensor in the field, including the four characterised in today's other instrument paper
(Werderman et al., \textit{Atmos Pollut Res}), is a light-scattering mass proxy with
essentially no response below $\sim$300\,nm. If Zhong et al.\ are right about the
regime, the distributed-sensing literature is building dense networks for a metric that
is becoming progressively less informative about the endpoint it is meant to proxy, and
the missing instrument --- a cheap, deployable surface-area or number-below-100\,nm
sensor --- is not on anyone's roadmap in this corpus.
\end{keybox}

\vspace{3mm}

% -------------------------------------------------------------- exec brief
\section*{\sffamily\bfseries\color{Deep}What changed today}
\vspace{-1.5mm}{\color{Amber}\rule{\linewidth}{1.1pt}}\vspace{2.5mm}

\footnotesize
\begin{itemize}
\item \textbf{The first four-way laboratory characterisation of the current low-cost
sensor generation --- and it arrived without an abstract.} Werderman et al.\ evaluate
\textbf{Sensirion SEN66, Piera IPS-7100, PurpleAir PIXEL and Alphasense OPC-N3} across
aerosol types and concentration ranges. Three of the four are units this watch has not
seen characterised before (the PIXEL is PurpleAir's move away from Plantower; the
IPS-7100 is one of the few consumer-priced units advertising particle-number output).
The publisher deposited title, authorship and DOI only, so it ships as
\textit{Metadata only} and is the highest-priority abstract retry on the queue. Last
author Sousan ran the 2021 OPC-N3/SPS30/AirBeam2/PMS-A003 evaluation that established the
environmental-vs-occupational calibration split, which is the lineage but not a finding.

\item \textbf{A biomonitoring result strong enough to be a sensor.} Giannetti et al.\
sampled \textit{Crematogaster scutellaris} at 30 traffic sites plus 4 controls in Parma
and found magnetic susceptibility of particles on the ant exoskeleton tracking traffic
intensity at \textbf{Spearman $\rho = 0.825$, $p < 0.001$}, with each additional
\textbf{1{,}000 vehicles/day associated with a 1.41-fold increase} and roughly
\textbf{700-fold across the sampled gradient}. Co-located soil showed nothing
($\rho = 0.167$, $p = 0.346$), which is the control that makes the result interesting:
the ant is integrating the airborne fraction, not the deposited one. SEM-EDS puts Fe-oxide
particles doped with Zn, Cu and Ti on the thorax, morphologically consistent with brake
and tyre wear. No calibration to $\mu$g/m$^3$ and no CI on the slope --- but as a
zero-power, zero-maintenance spatial sampler it is a genuinely different point in the
network-design space.

\item \textbf{Mediation of the air-pollution/cognition link is still not identified.}
Sakhvidi et al.\ find only \textbf{16 studies and 72 exposure--mediator--outcome
associations} across six pollutants for the whole literature, with cardio- and
cerebrovascular disease the most consistent candidate mediator and mental health, sleep
and lung function unresolved. Reporting quality was judged good; risk of bias was not,
particularly for outcome ascertainment and confounding. The honest reading is that the
mediation question is under-powered rather than answered.

\item \textbf{UK Biobank puts PM next to plasma neurodegeneration markers and finds the
interaction mostly absent.} In 19{,}645 participants aged $\geq$50 with complete
proteomic, metabolomic and environmental data, \textbf{GFAP was associated with dementia
at HR 2.65 (95\% CI 2.17--3.25)} and NfL with transitions to Parkinson's disease,
dementia and death. \textbf{PM $\times$ NfL interacted with mortality only at the nominal
level, and there was little evidence of PM interaction with either marker for dementia};
NO$_2$/NO$_x$ $\times$ GFAP was \emph{negative}. This is largely a null for the PM term
and should be cited that way, not as support for a PM--neurodegeneration pathway.

\item \textbf{No forest plot today.} The only CI-bearing estimate in the corpus has a
plasma protein on its left-hand side, not a particulate contrast, and Zhong's and
Giannetti's numbers carry no intervals. f5 is omitted --- the same call as 20 August.
\end{itemize}

\vspace{3mm}

% ==================================================================== ANALYTICS
\section*{\sffamily\bfseries\color{Deep}Corpus analytics}
\vspace{-1.5mm}{\color{Amber}\rule{\linewidth}{1.1pt}}\vspace{3mm}

\noindent\includegraphics[width=\linewidth]{\FIGDIR/f1_subtopics.png}
\figcap{Subtopic assignment of the 9 in-scope records; single-label by dominant
contribution. \textit{Sensing, forecasting \& instrumentation} and \textit{Exposure
assessment \& modelling} take 3 each --- \textbf{six of nine records sit on the
measurement side of the field}, the most method-weighted day this watch has produced.
\textbf{Only four of ten canonical clusters are occupied.} The six empty ones are
mechanistic toxicology, cardiovascular/metabolic, respiratory/allergic,
reproductive/developmental, occupational/indoor and other clinical endpoints --- that is,
the entire clinical half of the taxonomy is absent from a nine-record day.}

\vspace{2mm}
\noindent\includegraphics[width=\linewidth]{\FIGDIR/f2_design_endpoint.png}
\figcap{Left --- study architecture, after grouping. The distribution is unusually flat:
no group holds more than two of nine records --- measurement campaign 2 (the Parma ant
survey and the Fuzhou LDSA record), chamber/laboratory 2 (the CIMS calibration and the
biomass-burning chamber--aircraft work), \textit{Metadata only} 2, and one each of
observational cohort, review/synthesis and modelling/inventory. The abstract-free share
is 22\%, less than half of Saturday's 46\%. Right --- health endpoint by organ system.
Only four records carry a health endpoint at all (two neuro/cognitive, one respiratory,
one cancer); the \textit{No health endpoint} and \textit{Not extracted} bars together are
five of nine.}

\vspace{2mm}
\noindent\includegraphics[width=\linewidth]{\FIGDIR/f3_metric_geography.png}
\figcap{Left --- particle metric. Note that only one record is PM$_{2.5}$ mass in the
conventional sense; the rest use lung-deposited surface area, magnetic susceptibility of
deposited particles, aerosol oxidation state, OVOC precursor concentration, a composite
microclimate index, or GBD risk-factor categories. \textbf{The day's metric diversity is
itself the finding} and it is the direct visual counterpart to the LDSA/mass decoupling in
the lead item. Right --- geographic provenance. The \textit{global / multi-region} bar of 4
is not four global studies: it is two genuinely global records (the GBD burden analysis and
the cognitive-aging review), one two-country instrument campaign, and one metadata-only
record whose study city the publisher did not deposit. Europe 2 (Parma, UK), North America
2, China 1.}

\vspace{2mm}
\noindent\includegraphics[width=\linewidth]{\FIGDIR/f6_lifecourse.png}
\figcap{Life-course windows addressed, counted only where a record's study population is
explicitly anchored at that stage. Two of five windows occupied and both at the old end
--- the UK Biobank $\geq$50 cohort and the cognitive-aging review at 65+, the GBD female
lung-cancer burden across working-age adulthood. \textbf{No record this window studies a
population under 18}, and the perinatal window, occupied yesterday by the aerosol-acidity
preterm-birth stub, is empty again.}

\vspace{2mm}
\noindent\includegraphics[width=\linewidth]{\FIGDIR/f4_heatmap.png}
\figcap{Subtopic $\times$ study-architecture cross-tabulation. With nine records across
seven architectures the matrix is close to a permutation --- no cell holds more than one
record except \textit{Metadata only}, which spans sensing and exposure modelling. There is
no methodological concentration to report on a day this small, and reading a pattern into
it would be over-interpretation.}

\vspace{3mm}
% ==================================================================== DIGEST
\section*{\sffamily\bfseries\color{Deep}Paper-level digest}
\vspace{-1.5mm}{\color{Amber}\rule{\linewidth}{1.1pt}}\vspace{1mm}

{\fontsize{7.4}{9}\selectfont\color{Slate}Ordered by cluster. Each entry gives the
methodological core and the finding that matters, not an abstract paraphrase.\par}

\band{Sensing, forecasting \& instrumentation\hfill 3 records}{Amber}

\paperentry{Werderman et al.\ 2026}
{Laboratory evaluation of Sensirion SEN66, Piera IPS-7100, PurpleAir PIXEL and OPC-N3 low-cost aerosol sensors across aerosol types and concentration ranges}
{Atmos Pollut Res}
{\textbf{Metadata only.} Crossref by-ISSN sweep; title, authorship (Werderman, Richards,
Berkuta, Sousan) and DOI verified, no abstract deposited and no numbers extractable. It is
nevertheless the single most on-topic record of the window: a chamber characterisation of
\emph{four} current-generation units across aerosol type and concentration range, three of
which --- SEN66, IPS-7100, PurpleAir PIXEL --- have no prior evaluation in this corpus,
against the OPC-N3 as the established reference point. The same last author's 2021
four-sensor evaluation established that environmental and occupational settings need
separate calibration factors, so the plausible contribution here is the aerosol-type
dependence of that factor for the newer optics. \textbf{Highest-priority full-text
retrieval on the next run}; nothing is asserted about its findings until then.}
{10.1016/j.apr.2026.103185}

\paperentry{Giannetti et al.\ 2026}
{Ants as bioindicators of urban traffic pollution}
{Environ Pollut}
{Thirty traffic-graded sites plus four controls in Parma; magnetic susceptibility and
SEM-EDS on \textit{Crematogaster scutellaris} exoskeleton and on co-located soil at 5\,cm.
\textbf{Ant magnetic susceptibility tracked traffic intensity at $\rho = 0.825$
($p < 0.001$); soil did not ($\rho = 0.167$, $p = 0.346$)}. On a log scale the relation is
linear ($R = 0.66$, $p < 0.001$) with \textbf{1.41-fold per additional 1{,}000
vehicles/day}, $\sim$700-fold across the gradient. SEM-EDS found Fe-oxide particles doped
with Zn, Cu and Ti concentrated on the thorax, morphologically consistent with
traffic-derived ultrafine PM. The weaknesses are that no slope CI is given, magnetic
susceptibility is a proxy for the ferrous fraction only (it says nothing about organic or
secondary mass), and nest-density confounding with traffic corridors is not addressed. For
distributed sensing the interest is orthogonal: a passive, unpowered integrator with
month-scale averaging at effectively zero unit cost, usable as a siting prior for where to
put real instruments rather than as a replacement for them.}
{10.1016/j.envpol.2026.129010}

\paperentry{Roska et al.\ 2026}
{Beyond targeted calibrations: comprehensive quantification of oxygenated organic compounds in urban atmospheres}
{Environ Sci Technol}
{A collision-induced-dissociation voltage-scanning method to derive sensitivities for an
ammonium-adduct CIMS, so OVOCs without an authentic standard can still be quantified.
Validated in the laboratory, then applied to the \textbf{2023 AEROMMA airborne campaign
and the 2024 SAPHIR-CHANEL chamber campaign}: coverage rose to \textbf{98\% of total
signal in chamber and 78\% in the field}, and adding the newly quantified species raised
estimated \textbf{OH reactivity by $\sim$5\% and SOA formation potential by $\sim$62\%}.
Coverage degrades with molecular size and oxidation state, and chamber-derived
sensitivities are being transferred to ambient air on the strength of trend agreement
rather than field standards --- the paper says so. The 62\% figure is the one that matters
downstream: SOA yields used in chemical-transport models that feed exposure surfaces are
calibrated against measurements that were missing most of this signal.}
{10.1021/acs.est.6c04149}

\band{Exposure assessment \& modelling\hfill 3 records}{Deep}

\paperentry{Zhong et al.\ 2026}
{Photochemical ultrafine particles pose potential inflammatory risk despite low PM$_{2.5}$ mass}
{Environ Int}
{Three years of hourly Fuzhou observations; lung-deposited surface area compared against
what the local LDSA/PM$_{2.5}$ ratio predicts, with predicted inflammation from an
animal-derived surface-area dose--response. \textbf{6.6\% of PM$_{2.5}$-attainment hours
($<$35\,$\mu$g/m$^3$) carried 7.5$\times$ the expected surface-area exposure, and 82\% of
the predicted inflammatory risk in those hours was undetected by mass}. Median
PM$_{2.5}$ during them was \textbf{3.0\,$\mu$g/m$^3$}, with the dose--response projecting
27\% lung neutrophil infiltration. Machine learning attributes the events to photochemical
production of 50--100\,nm particles on warm, high-radiation summer--autumn afternoons
(O$_3$ and secondary organic carbon dominate importance). \textbf{The limitation that
threatens the inference is the dose--response itself}: an animal-derived surface-area
relation extrapolated to a human population gives a predicted, not an observed,
inflammatory response, and no human endpoint is measured here. Even so, the decoupling of
LDSA from mass is measured, not modelled, and it is a direct challenge to the assumption
that a dense network of scattering-based mass sensors indexes health-relevant exposure.}
{10.1016/j.envint.2026.110469}

\paperentry{Zhang et al.\ 2026b}
{Photoinduced intraparticle chemistry drives the rapid atmospheric oxidation of biomass burning organic aerosol}
{Environ Sci Technol}
{Chamber isolation of biomass-burning aerosol aging pathways --- UV-A irradiation,
OH-driven SOA formation, heterogeneous OH oxidation, ozonolysis and evaporation ---
benchmarked against aircraft measurements of wildfire and prescribed-burn plumes.
\textbf{UV-A irradiation drove intraparticle aging an order of magnitude faster than
gas-phase oxidation}, at an estimated quantum yield of $\sim$0.4\% (four in a thousand
absorbed photons generate particle-phase radicals). The authors estimate that
\textbf{over half of oxygenated OA in biomass-burning-influenced regions may originate
from photo-oxidised primary OA within the first day}. This inverts the standing
assumption that gas-phase oxidation dominates BBOA aging, and it matters for exposure
work because oxidation state governs both the oxidative potential of the inhaled particle
and the optical properties that low-cost scattering sensors respond to --- a smoke plume's
sensor bias is not constant along its trajectory.}
{10.1021/acs.est.6c02532}

\paperentry{Mohammed 2026}
{Spatiotemporal dynamics of atmospheric vulnerability in hot arid cities: introducing the Urban Microclimate Air Stress Index (UMASI)}
{Build Environ}
{\textbf{Metadata only.} Crossref sweep, no abstract deposited, no study city named in the
metadata. A composite index combining microclimate and air-quality stress is exactly the
kind of construct that needs its weighting scheme inspected before it is used, so nothing
is characterised here. Flagged for retry; low priority relative to Werderman et al.}
{10.1016/j.buildenv.2026.115166}

\clearpage

\band{Neuro / mental health\hfill 2 records}{Deep}

\paperentry{Beydoun et al.\ 2026}
{The interplay of air pollution with plasma neurodegenerative markers and metabolome for dementia, Parkinson's disease and all-cause mortality}
{Ecotoxicol Environ Saf}
{UK Biobank, \textbf{19{,}645 participants aged $\geq$50} with complete proteomic,
metabolomic and environmental data; Cox models plus multistate Weibull models for
transitions from healthy to dementia, Parkinson's disease and death. Higher NfL predicted
all three transitions; \textbf{GFAP was specific to dementia, HR 2.65 (95\% CI
2.17--3.25)}. On the exposure side the result is largely negative: \textbf{PM $\times$
NfL interacted with mortality only at the nominal level, NO$_2$/NO$_x$ $\times$ GFAP
interacted \emph{negatively}, and there was little evidence of PM interaction with either
marker for dementia}. Metabolomic principal components (lipid, amino-acid, energy) did
modify the biomarker--outcome relations and interacted with PM for PD and dementia; a
branched-chain-amino-acid component attenuated the GFAP--dementia association. The
interaction tests are nominal and unadjusted for multiplicity across a large grid, so the
positive ones should be treated as hypothesis-generating. Cite this for the biomarker
result and for the \emph{absence} of a PM interaction, not for a PM effect.}
{10.1016/j.ecoenv.2026.120657}

\paperentry{Sakhvidi et al.\ 2026}
{Ambient air pollution exposure and cognitive aging; mediators and possible pathways: a systematic review on the epidemiological studies}
{Ageing Res Rev}
{\textbf{16 studies, 72 unique exposure--mediator--outcome associations}, six pollutants
(PM$_{2.5}$, PM$_{10}$, NO$_2$, NO$_x$, black carbon, PM$_1$ components), with PM$_{2.5}$
in 65\% of analyses. Outcomes span memory, processing speed and diagnosed dementia;
candidate mediators span depression/anxiety/stress, FEV$_1$/FVC/PEF, stroke and
hypertension, CRP, type 2 diabetes, sleep and neuroanatomical change. \textbf{Cardio- and
cerebrovascular disease emerged as the most consistent mediator}; mental health, sleep and
lung function showed potential but unconfirmed mediation. Reporting quality was judged
generally good, but most studies carried some risk of bias in outcome assessment and
confounding adjustment. No pooled estimate, and 16 studies across six pollutants and this
many candidate mediators is a thin base --- the useful output is the map of which pathways
have been tested at all, and the observation that almost none have been tested with a
formal causal-mediation design.}
{10.1016/j.arr.2026.103330}

\band{Burden, policy \& mitigation\hfill 1 record}{Coral}

\paperentry{Wang et al.\ 2026}
{Trends, inequalities, and 2030 projections of women's tracheal, bronchial, and lung cancer attributable to air pollution, radon, and tobacco exposure}
{Front Glob Womens Health}
{GBD 2023, female tracheal/bronchial/lung cancer 1990--2023, five inhalation-related risk
factors, with SDI-gradient smoothing, concentration indices for inequality, decomposition
by demographic vs epidemiological change, and projection to 2030. Smoking remains the
leading contributor, \textbf{ambient particulate matter second, secondhand smoke third}.
Over 1990--2023 \textbf{household solid-fuel-attributable mortality fell most sharply
while particulate-matter-attributable mortality rose slightly}, with China the standout
for the PM-attributable share. Along the SDI gradient, PM shows a \textbf{rising burden
with diminishing inequality} --- the opposite pattern to radon and tobacco --- and
decomposition attributes the PM increase to population growth and aging rather than to
rising risk. Projected to keep rising to 2030. Standard GBD caveats bind hard here:
the exposure is a modelled global surface, the attributable fraction rests on assumed
counterfactual and RR functions, and the 2030 projection is an extrapolation of fitted
trend, not a scenario.}
{10.3389/fgwh.2026.1662451}

\vspace{4mm}
% ==================================================================== TABLE
\section*{\sffamily\bfseries\color{Deep}Corpus register}
\vspace{-1.5mm}{\color{Amber}\rule{\linewidth}{1.1pt}}\vspace{2.5mm}

{\fontsize{6.7}{8.4}\selectfont
\setlength{\tabcolsep}{2.4pt}
\begin{longtable}{@{}>{\RaggedRight}p{25mm}>{\RaggedRight}p{41mm}>{\RaggedRight}p{27mm}%
>{\RaggedRight}p{23mm}>{\RaggedRight}p{18mm}>{\RaggedRight}p{15mm}@{}}
\toprule
\textsf{\textbf{First author}} & \textsf{\textbf{Focus}} & \textsf{\textbf{Journal}} &
\textsf{\textbf{Design}} & \textsf{\textbf{Metric}} & \textsf{\textbf{Region}}\\
\midrule
\endfirsthead
\toprule
\textsf{\textbf{First author}} & \textsf{\textbf{Focus}} & \textsf{\textbf{Journal}} &
\textsf{\textbf{Design}} & \textsf{\textbf{Metric}} & \textsf{\textbf{Region}}\\
\midrule
\endhead
\multicolumn{6}{@{}l@{}}{\cellcolor{Mist}\textsf{\textbf{\textcolor{Deep}{Neuro / mental health}}}}\\[0.6mm]
\href{https://doi.org/10.1001/jamanetworkopen.2026.25015}{Dales} & Wildfire-sourced PM$_{2.5}$ and ED visits for migraine and other primary headache & \textit{JAMA Netw Open} & Case-crossover & PM$_{2.5}$ (wildfire) & Canada\\
\href{https://doi.org/10.1093/ageing/afag215}{Glover} & Umbrella meta-review of urban environmental exposures and cognitive health & \textit{Age Ageing} & Umbrella review / meta-analysis & PM$_{2.5}$ & Global\\
\href{https://doi.org/10.1016/j.nsa.2026.107019}{Goossens} & In utero air pollution exposure and pre-/postnatal brain development & \textit{Neurosci Appl} & Narrative review & PM$_{2.5}$ / PM$_{10}$ / PM$_{\mathrm{coarse}}$ & Global\\
\href{https://doi.org/10.1016/j.isci.2026.116617}{Ma} & PM$_{2.5}$-induced neurotoxicity and depression: exposure risks and mechanisms & \textit{iScience} & Narrative review & PM$_{2.5}$ + components & Global\\
\href{https://doi.org/10.3760/cma.j.cn112150-20250804-00752}{Pan} & Maternal PM exposure and offspring neurodevelopment: effects and mechanisms & \textit{Zhonghua Yu Fang Yi Xue Za Zhi} & Narrative review & PM (bulk) & China\\
\addlinespace[0.9mm]
\multicolumn{6}{@{}l@{}}{\cellcolor{Mist}\textsf{\textbf{\textcolor{Teal}{Cardiovascular \& metabolic}}}}\\[0.6mm]
\href{https://doi.org/10.3760/cma.j.cn112338-20251120-00843}{Huang} & PM$_{2.5}$ and dyslipidaemia in rural Chinese children; dietary effect modification & \textit{Zhonghua Liu Xing Bing Xue Za Zhi} & Cross-sectional (multistage) & PM$_{2.5}$ & China\\
\href{https://doi.org/10.32394/pe/218966}{Kazmierska} & Air pollution and cardiovascular disease: review of current data & \textit{Przegl Epidemiol} & Narrative review & PM$_{2.5}$ / PM$_{10}$ & Global\\
\href{https://doi.org/10.1007/s00125-026-06814-2}{Xi} & Air pollution, built environment and temperature vs. prediabetes phenotypes & \textit{Diabetologia} & Prospective cohort & PM$_{2.5}$ & Germany\\
\href{https://doi.org/10.1161/JAHA.125.046917}{Yang} & Short-term PM$_{2.5}$ and hospitalisation/ED use in atrial fibrillation & \textit{J Am Heart Assoc} & Case-crossover & PM$_{2.5}$ (wildfire-influenced) & USA\\
\addlinespace[0.9mm]
\multicolumn{6}{@{}l@{}}{\cellcolor{Mist}\textsf{\textbf{\textcolor{Sky}{Reproductive \& developmental}}}}\\[0.6mm]
\href{https://doi.org/10.1016/j.jes.2025.10.043}{Chen} & Prenatal air pollution, fetal growth trajectory phenotypes and child neurodevelopment & \textit{J Environ Sci (China)} & Prospective cohort & PM$_{2.5}$ / PM$_{10}$ & China\\
\href{https://doi.org/10.1016/j.ecoenv.2026.120566}{Huang} & Joint prenatal bisphenol + air-pollutant mixture exposure and preschool IQ & \textit{Ecotoxicol Environ Saf} & Prospective cohort (mixtures) & PM$_{2.5}$ / PM$_{10}$ & China\\
\href{https://doi.org/10.1007/s40572-026-00552-8}{Peterson} & Ambient air pollution, PAHs and pubertal development: critical review & \textit{Curr Environ Health Rep} & Critical review & PM$_{2.5}$ components & Global\\
\href{https://doi.org/10.1016/j.envres.2026.125313}{Xie} & Long-term PM$_{2.5}$ and timing of pubertal development in boys and girls & \textit{Environ Res} & Prospective cohort & PM$_{2.5}$ & China\\
\addlinespace[0.9mm]
\multicolumn{6}{@{}l@{}}{\cellcolor{Mist}\textsf{\textbf{\textcolor{Sage}{Respiratory \& allergic}}}}\\[0.6mm]
\href{https://doi.org/10.1007/s00484-026-03289-z}{Bayar Muluk} & Climate change and otolaryngology: airway disease impacts and emission reduction & \textit{Int J Biometeorol} & Narrative review & PM$_{2.5}$ & Global\\
\href{https://doi.org/10.4103/mgr.MEDGASRES-D-25-00164}{Liao} & Short-term air pollution, respiratory symptoms and lung function in asthma & \textit{Med Gas Res} & Panel study & PM$_{2.5}$ / PM$_{10}$ & China\\
\href{https://doi.org/10.3390/epidemiologia7040097}{Sanchez de Prada} & Ambient PM and asthma-coded admissions in children 0-3 y, Spain & \textit{Epidemiologia} & Ecological panel & PM$_{2.5}$ / PM$_{10}$ & Spain\\
\href{https://doi.org/10.1016/j.jhazmat.2026.143076}{Yang} & Allergenic vs. toxic risk of airborne seed hairs compared with ambient PM$_{2.5}$ & \textit{J Hazard Mater} & Field measurement & PM$_{2.5}$ + bioaerosol & China\\
\addlinespace[0.9mm]
\multicolumn{6}{@{}l@{}}{\cellcolor{Mist}\textsf{\textbf{\textcolor{Violet}{Mechanistic toxicology}}}}\\[0.6mm]
\href{https://doi.org/10.1016/j.jes.2025.12.074}{Ding} & PM$_{2.5}$ promotes vascular calcification via miR-27b-5p/HES1 signalling & \textit{J Environ Sci (China)} & In vitro + organoid + murine & PM$_{2.5}$ & China\\
\href{https://doi.org/10.1007/s00360-026-01698-5}{Maraschi} & Osmoregulatory disruption and metal bioaccumulation from settleable atmospheric PM & \textit{J Comp Physiol B} & Controlled exposure (ecotox) & Settleable APM & Brazil\\
\href{https://doi.org/10.1016/j.jes.2025.11.018}{Wu} & Macrophage NLRP3 inflammasome in real-ambient PM-induced hepatic glucose dysregulation & \textit{J Environ Sci (China)} & Real-ambient murine + co-culture & Real-ambient PM & China\\
\addlinespace[0.9mm]
\multicolumn{6}{@{}l@{}}{\cellcolor{Mist}\textsf{\textbf{\textcolor{Clay}{Occupational \& indoor}}}}\\[0.6mm]
\href{https://doi.org/10.3389/fmedt.2026.1796828}{Chen} & Nested suction instrument to cut intraoperative surgical-smoke PM$_{2.5}$ & \textit{Front Med Technol} & Bench / device experiment & PM$_{2.5}$ (surgical smoke) & China\\
\href{https://doi.org/10.1093/annweh/wxag055}{Cole} & Determinants of respirable crystalline silica exposure in Australian tunnelling & \textit{Ann Work Expo Health} & Observational exposure study & Respirable dust / RCS & Australia\\
\href{https://doi.org/10.1016/j.jhazmat.2026.143081}{Ding} & Stage-resolved hair Zn-isotope dynamics as a firefighting combustion-exposure tracer & \textit{J Hazard Mater} & Repeated-measures biomarker & Combustion-derived PM & China\\
\href{https://doi.org/10.3390/jox16040119}{Ramsden} & Xenobiotic hazards in aircraft cabin air & \textit{J Xenobiot} & Review + risk economics & Ultrafine particles & Global\\
\addlinespace[0.9mm]
\multicolumn{6}{@{}l@{}}{\cellcolor{Mist}\textsf{\textbf{\textcolor{Amber}{Sensing, forecasting \& instrumentation}}}}\\[0.6mm]
\href{https://doi.org/10.1016/j.jenvman.2026.130515}{Shi} & Adaptive spatiotemporal graph transformer for multi-site multi-step PM$_{2.5}$ forecasting & \textit{J Environ Manage} & Deep-learning model & PM$_{2.5}$ & China / India\\
\href{https://doi.org/10.3390/bios16070380}{Tsuchiya} & Seasonal and indoor-outdoor characteristics of fluorescent aerosol particles in offices & \textit{Biosensors} & Real-time sensor campaign & Bioaerosol / coarse & Japan\\
\addlinespace[0.9mm]
\multicolumn{6}{@{}l@{}}{\cellcolor{Mist}\textsf{\textbf{\textcolor{Amber}{Exposure assessment \& modelling}}}}\\[0.6mm]
\href{https://doi.org/10.1029/2025GH001621}{Solorzano-Araque} & Multiscale exposure modelling: LOTOS-EUROS + mobility surveys + particle toxicity & \textit{GeoHealth} & CTM + mobility integration & PM$_{2.5}$ / PM$_{10}$ & Colombia\\
\addlinespace[0.9mm]
\multicolumn{6}{@{}l@{}}{\cellcolor{Mist}\textsf{\textbf{\textcolor{Coral}{Burden, policy \& mitigation}}}}\\[0.6mm]
\href{https://doi.org/10.1002/hsr2.72886}{Karamzad} & Ischaemic heart disease burden in MENA 1990-2021; PM attributable fraction 28.2\% & \textit{Health Sci Rep} & GBD secondary analysis & Ambient PM & MENA\\
\href{https://doi.org/10.1371/journal.pone.0347989}{Rees} & HEPA filtration to reduce respiratory infection in care homes & \textit{PLoS One} & RCT process evaluation & Filtered PM & UK\\
\href{https://doi.org/10.1016/j.jes.2026.01.009}{Shan} & Policy-driven emission reductions and post-pandemic rebound, high-resolution inventories & \textit{J Environ Sci (China)} & Emission inventory & PM emissions & China\\
\addlinespace[0.9mm]
\multicolumn{6}{@{}l@{}}{\cellcolor{Mist}\textsf{\textbf{\textcolor{Slate}{Other clinical endpoints}}}}\\[0.6mm]
\href{https://doi.org/10.3390/dj14070440}{Lupita} & PM$_{2.5}$, preventive-care capacity and caries/edentulism burden across the EU & \textit{Dent J} & Ecological panel & PM$_{2.5}$ & EU-27\\
\href{https://doi.org/10.1016/j.jdin.2026.06.008}{Park} & Long-term PM exposure and risk of androgenetic alopecia & \textit{JAAD Int} & Nationwide cohort & PM$_{2.5}$ / PM$_{10}$ & South Korea\\
\href{https://doi.org/10.3171/2026.3.PEDS2615}{Sun} & Environmental exposure and area deprivation vs. paediatric hydrocephalus outcomes & \textit{J Neurosurg Pediatr} & Registry cohort + geospatial & PM$_{2.5}$ (satellite-derived) & USA\\
\addlinespace[0.9mm]
\bottomrule
\end{longtable}}


\vspace{4mm}
\begin{keybox}
{\sffamily\bfseries\footnotesize\color{Deep}Provenance and method}\\[1.2mm]
{\fontsize{7.2}{8.8}\selectfont
Entry window \textbf{2026-08-23 $\rightarrow$ 2026-08-23}, contiguous with the 22 August
issue. \textbf{This issue was built late}, on 24 August, because the scheduled Sunday run
did not fire; the harvest window is unchanged and no gap was left, but the record is noted
rather than concealed. \textbf{Sources.} PubMed E-utilities on \texttt{[EDAT]} returned
\textbf{1 record} (a ruminant-methane study, screened out) --- the lowest single-day
PubMed yield the watch has recorded, consistent with a Sunday deposit floor; the
connector re-query on both the health and the instrumentation axis returned zero for the
same window. Europe PMC on \texttt{CREATION\_DATE} returned 20 records, of which
\textbf{13 were archival PMC back-scans} (publication years 1974--2005, no title in the
record) and 5 shipped. The Crossref by-ISSN sweep returned 12 records after front-matter
filtering and shipped 4. \textbf{OpenAlex returned HTTP 429 for the 21st consecutive
run.} arXiv returned nothing new under the 3-day recency screen. Consensus and Scholar
Gateway were both queried for the Werderman et al.\ sensor evaluation and neither has
indexed it yet, which is why it ships as a stub.
\textbf{Screening.} 22 records were screened out and written to
\texttt{state/rejected.jsonl} with an individual reason each; the 13 archival back-scans
carry a shared structural reason. One borderline exclusion is logged as borderline:
a national wastewater-based-epidemiology assessment, environmental and
population-scale but with no airborne route.
\textbf{Labels.} Two records carry no abstract and are marked \textit{Metadata only};
their design, sample size, endpoint and study location are recorded as not extracted
rather than inferred. \textbf{No forest plot} is included --- see the closing bullet of
``What changed today'' for why. All DOIs were checked by \texttt{check\_dois.py} before
publication. Content derived from publisher metadata; abstracts remain the copyright of
their respective publishers.\par}
\end{keybox}

\end{document}
